\chapter{Core Concepts}
\label{chap:core_concepts}

This chapter introduces the actuarial concepts that underpin \actsim. Users
familiar with frequency--severity modelling and stochastic simulation may
skim this material, but the notation introduced here is used throughout the
manual.

\section{Frequency Distributions}
\label{sec:frequency_distributions}

The \emph{frequency} of claims is the number of claims that occur during a
specified period (often a policy year). \actsim models claim frequency
using discrete probability distributions. The most common choices are:

\begin{itemize}[leftmargin=*]
    \item \textbf{Poisson distribution.} A simple one-parameter distribution
          where the mean equals the variance. The probability mass
          function is
          \[
              P(N = n) = \frac{\lambda^{n} e^{-\lambda}}{n!},
              \qquad n = 0, 1, 2, \dots
          \]
          where $\lambda > 0$ is the expected number of claims.
    \item \textbf{Negative binomial distribution.} A two-parameter
          distribution that allows variance greater than the mean
          (over-dispersion), often more realistic for claim counts.
\end{itemize}

\section{Severity Distributions}
\label{sec:severity_distributions}

The \emph{severity} of a claim is the loss amount associated with that
claim. \actsim supports a wide range of continuous severity distributions,
including:

\begin{itemize}[leftmargin=*]
    \item Lognormal
    \item Gamma
    \item Weibull
    \item Pareto
    \item Exponential
    \item Beta
    \item Normal, logistic, uniform (occasionally used for sensitivity
          analysis or stylised examples)
\end{itemize}

For a continuous severity $X$ with density $f_X$ and cumulative
distribution function $F_X$, the conventional moments are
\[
    E[X] = \int_0^{\infty} x f_X(x)\, dx, \qquad
    \mathrm{Var}(X) = E[X^2] - (E[X])^{2}.
\]

\section{Aggregate Loss}
\label{sec:aggregate_loss}

The \emph{aggregate loss} for a period is the total loss across all claims
that occur during that period. If $N$ denotes the number of claims and
$X_1, X_2, \dots, X_N$ denote the individual claim severities, then
the aggregate loss is
\[
    S = \sum_{i=1}^{N} X_i.
\]

When $N$ and the $X_i$ are mutually independent and the $X_i$ are
independent and identically distributed, $S$ is referred to as a
\emph{compound distribution}. Standard moment results in this setting
are
\[
    E[S] = E[N] \cdot E[X], \qquad
    \mathrm{Var}(S) = E[N] \cdot \mathrm{Var}(X) + \mathrm{Var}(N) \cdot (E[X])^{2}.
\]

In \actsim, aggregate losses are typically estimated by simulation rather
than by analytical closed-form expressions, which makes the framework
flexible across distribution choices.

\section{Dependency and Correlation}
\label{sec:dependency}

In practice, claim frequency and severity are not always independent. For
example, a year with many claims may also feature larger average severities
due to a common driver. \actsim supports two mechanisms for introducing
dependency between frequency and severity:

\begin{enumerate}[leftmargin=*]
    \item \textbf{Linear correlation.} A user-supplied correlation
          coefficient is applied to a bivariate normal copula structure
          via Cholesky decomposition.
    \item \textbf{Copula-based dependency.} \actsim supports several
          standard copulas, including Gaussian, Frank, Gumbel, and
          Clayton copulas, parametrised by a dependency parameter.
\end{enumerate}

Multivariate dependency across multiple lines of business is also
supported through correlated marginal simulation, configured via a
correlation matrix and a list of marginal distributions.

\section{Claim-Level Simulation}
\label{sec:claim_level}

Many actuarial workflows require not just aggregate losses but also
\emph{claim-level} data: individual claim records with policy identifiers,
occurrence dates, severities, and development histories. \actsim's
\texttt{ClaimSimulator} class generates these records by:

\begin{enumerate}[leftmargin=*]
    \item Grouping policies by their frequency and severity distribution
          specifications.
    \item Simulating claim counts and severities for each group using the
          underlying \texttt{StochasticSimulator}.
    \item Assigning occurrence dates within each policy's exposure window
          using a non-homogeneous Poisson process.
    \item Optionally applying claim development to each simulated claim.
\end{enumerate}

\section{Claim Development}
\label{sec:claim_development_concept}

Claims do not typically reach their ultimate cost on the date they occur.
Reported and paid amounts develop over time as additional information
becomes available. \actsim represents this development through
user-supplied \emph{loss development factors} (LDFs).

If $U$ denotes the ultimate loss for a claim and $\mathrm{CDF}_a$ denotes
the cumulative development factor to ultimate at age $a$, then the
reported loss at age $a$ is approximated by
\[
    L_a = \frac{U}{\mathrm{CDF}_a}.
\]

The incremental development between two adjacent ages is
\[
    \Delta L_a = L_a - L_{a-1}.
\]

Stochastic variation around expected LDFs is introduced via a normal
multiplier with user-specified volatility.

\section{Risk Metrics}
\label{sec:risk_metrics_concept}

\actsim provides several standard risk metrics, including:

\begin{itemize}[leftmargin=*]
    \item \textbf{Mean and standard deviation} of aggregate losses.
    \item \textbf{Percentiles} of the loss distribution.
    \item \textbf{Value at Risk (VaR)} at user-specified confidence levels.
    \item \textbf{Tail Value at Risk (TVaR)}, also known as conditional
          tail expectation.
    \item \textbf{Aggregate Exceedance Probability (AEP)}, related to
          aggregate (annual) losses.
    \item \textbf{Occurrence Exceedance Probability (OEP)}, related to
          maximum single-event losses within a year.
\end{itemize}

These metrics are defined formally in Chapter~\ref{chap:risk_metrics}.
